\label{sec:results}

\subsection{Purpose of the Court Usage Template}

The court usage template in this section is the primary contribution of M.8.
Individual M-track papers (M.1--M.7) each provide a single-signal expert
witness template.
M.8 provides the synthesis template: the complete statement an expert witness
should make when the composite community character index was used to draw
the map.

The template is written for use in:
\begin{itemize}
  \item Congressional redistricting proceedings before federal or state courts
  \item State legislative redistricting proceedings before state courts or
    redistricting commissions
  \item Special master proceedings in cases where a court has assumed
    redistricting authority
  \item Commission public comment and expert report contexts where quantitative
    community-of-interest evidence is solicited
\end{itemize}

The template has five sections: data provenance, community index description,
non-partisan declaration, specific community identification, and algorithmic
transparency.

\subsection{Expert Witness Statement Template}

\subsubsection*{Section 1: Data Provenance}

\begin{quote}
``The submitted congressional plan was drawn using the \textsc{bisect}
algorithmic redistricting platform, which is a public-domain computational
tool for graph-partitioning-based congressional district design.
The plan was produced by optimizing the Composite Community Character Index
(CCCI), a quantitative measure of community bond strength derived from
seven federal government administrative and remote sensing datasets.
No electoral, racial, or partisan data entered any component of the index.

The seven data sources are:
\begin{enumerate}
  \item U.S.\ Census Bureau TIGER/Line shapefiles (school district, municipal,
    and fire district boundaries), 2020.
    Source: \url{https://www.census.gov/geographies/mapping-files/time-series/geo/tiger-line-file.html}.
  \item U.S.\ Census Bureau LEHD LODES Workplace Area Characteristics
    (WAC), version~8, 2020.
    Source: \url{https://lehd.ces.census.gov/data/}.
  \item USGS National Land Cover Database (NLCD) 2021, 30-meter resolution.
    Source: \url{https://www.mrlc.gov/data/nlcd-2021-land-cover-conus}.
  \item U.S.\ Census Bureau American Community Survey 5-year estimates,
    2016--2020 (housing tenure, age, value, and renter fraction by tract).
    Source: \url{https://www.census.gov/programs-surveys/acs/}.
  \item U.S.\ Census Bureau LEHD LODES Origin-Destination file (OD),
    version~8, 2020.
    Source: \url{https://lehd.ces.census.gov/data/}.
  \item USGS Shuttle Radar Topography Mission (SRTM) 30-meter elevation data.
    Source: \url{https://earthexplorer.usgs.gov/}.
  \item GTFS (General Transit Feed Specification) public transit schedules
    for applicable metropolitan areas, 2020 vintage.
    Source: TransitFeeds.com and individual transit authority GTFS publications.
\end{enumerate}

All datasets are public domain or publicly licensed federal government products.
None contains voter registration data, election results, precinct boundaries,
or any other information about the political preferences of the populations
within the census tracts.''
\end{quote}

\subsubsection*{Section 2: Community Index Description}

\begin{quote}
``The Composite Community Character Index assigns a numerical score to every
pair of adjacent census tracts reflecting the strength of the community bond
between those tracts.
The score is a weighted average of seven component similarity measures,
each derived from one of the federal data sources listed above.

The seven components and their default weights are:
(1) zone co-membership (school district, municipal boundary, fire district),
weight 30\%; (2) economic character (commercial intensity, industrial fraction,
jobs per resident), weight 20\%; (3) land use similarity (residential,
commercial, open, and natural fractions from NLCD), weight 20\%;
(4) housing character (tenure, age, value, renter fraction), weight 15\%;
(5) commuting shed membership (weighted Jaccard similarity of LODES OD
destination distributions), weight 10\%;
(6) topographic community (terrain profile similarity from SRTM), weight 3\%;
(7) transit accessibility (GTFS-based transit service similarity), weight 2\%.

The weights sum to 100\%.
When data for a component is unavailable for a given state or tract pair,
that component is omitted from the average and the remaining weights are
renormalized.
The minimum signal set required for a valid composite run is zone
co-membership, economic character, and land use.

Adjacent census tracts with a high composite score (close to 1.0) are
strongly bonded communities on multiple dimensions.
Adjacent census tracts with a low composite score (close to 0.0) are
communities that differ substantially on multiple dimensions.
The \textsc{bisect} algorithm assigns district boundaries preferentially
to low-composite-score adjacencies---the natural seams between community
types---and keeps high-composite-score adjacencies within the same district.''
\end{quote}

\subsubsection*{Section 3: Non-Partisan Declaration}

\begin{quote}
``No electoral, racial, or partisan data entered any component signal of
the Composite Community Character Index.
The following data types were considered and explicitly excluded from
the computation:
\begin{itemize}
  \item Voter registration records (state or county election authority data)
  \item Election results (any precinct, county, or district level)
  \item Precinct boundary shapefiles used as a proxy for electoral geography
  \item Racial or ethnic composition data (decennial census P1/P2 tables or
    ACS race tables)
  \item Partisan lean indices (Cook, Sabato, or any other partisan rating)
\end{itemize}

Any correlation between the submitted plan and partisan outcomes is a consequence
of the geographic correlation between community character signals (employment
concentration, land use type, administrative zone membership) and partisan
sorting patterns in the American electorate.
That correlation is a feature of the population landscape, not a product of
the algorithm.
The algorithm did not have access to, and did not use, any information about
partisan voting patterns.

This satisfies the standard articulated in \textit{Rucho v.\ Common Cause},
588 U.S.\ 684 (2019), that federal courts evaluate partisan gerrymandering
by asking whether partisan considerations predominated in the drawing of
district lines.
No partisan consideration entered any component of the index.''
\end{quote}

\subsubsection*{Section 4: Specific Community Identification}

\begin{quote}
``The following communities of interest are preserved in the submitted plan
as demonstrated by the Composite Community Character Index:

[PRACTITIONER NOTE: Insert specific communities from the state's CCCI analysis.
Use the tract-level CCCI output to identify adjacency pairs with high composite
scores that are kept within the same district. For each preserved community,
cite at least two component signals. Example language:]

`The [community name] community in [county/counties] is preserved as a cohesive
unit within [district number].
This community has a composite score above 0.75 on the following dimensions:
zone co-membership (tracts in [community] share the [school district name]
school district and the [municipality] municipal service boundary, zone
co-membership score 0.95); economic character (tracts in [community] share
a [residential/commercial/industrial] employment profile, economic character
score 0.88); and land use similarity (tracts in [community] are predominantly
[residential/commercial/agricultural], land use score 0.82).
The composite score of 0.84 reflects a strongly bonded community on three
independent dimensions.'

[For water-separated communities using the M.2 hard-cut rule:]

`The [Outer Banks / barrier island / coastal peninsula] community in [county]
is separated from mainland tracts by a zero-weight water boundary.
Both the [community tracts] and the adjacent mainland tracts have majority
water classification under NLCD (open water fraction exceeds 50\%), triggering
the land use hard-cut rule.
This boundary corresponds to [body of water name], which this commission and
its predecessors have consistently recognized as a community-defining natural
boundary.'
''
\end{quote}

\subsubsection*{Section 5: Algorithmic Transparency}

\begin{quote}
``The \textsc{bisect} platform computes district assignments by solving the
following optimization problem: partition the census tract adjacency graph
into $k$ connected components of approximately equal population, minimizing
the total weight of edges that are cut.
The composite community character weights modify the edge weights so that
cuts are cheaper at community seams and more expensive within communities.

The specific formulas are:
\begin{align*}
  w_{\mathrm{comp}}(u, v) &=
    w_{\mathrm{geo}}(u, v) \times
    \frac{\sum_{i \in \mathcal{A}} \omega_i \cdot \mathrm{sim}_i(u, v)}
         {\sum_{i \in \mathcal{A}} \omega_i}
\end{align*}
where $w_{\mathrm{geo}}$ is the geographic boundary length between tracts
$u$ and $v$, $\omega_i$ are the component weights (zone 0.30, economic 0.20,
land 0.20, housing 0.15, commute 0.10, topographic 0.03, transit 0.02),
$\mathcal{A}$ is the set of available signals for this edge, and
$\mathrm{sim}_i(u, v)$ is the component similarity score for signal $i$.

The formula is symmetric ($w(u,v) = w(v,u)$), non-negative, and bounded
above by $w_{\mathrm{geo}}(u, v)$.
It has been implemented in the \textsc{bisect} Rust codebase and verified
by a suite of L0 unit tests including invariants for non-negativity,
symmetry, and graceful degradation.

The graph partitioning is performed by the METIS library
\citep{karypis1998fast} using the minimum-cut objective with population-balance
constraints.
METIS is a well-established, peer-reviewed graph partitioning library used
in high-performance computing, finite element analysis, and computational
science.
Its application to redistricting is a natural extension of its graph
partitioning capabilities.

The population balance constraint is set at $\pm$[X]\% of ideal population,
where the ideal population is the state's total population divided by $k$.
This constraint is enforced by METIS during the partitioning step.
Equal population is the first criterion; community character weights apply
only within the space of legally valid equal-population plans.''
\end{quote}

\subsection{Brief Citation Template}

For use in redistricting briefs, the following formulation provides a concise
reference to the composite index analysis:

\begin{quote}
``The proposed congressional plan was produced using an algorithmic tool
(the \textsc{bisect} platform) that operationalizes communities of interest
through the Composite Community Character Index, a weighted average of seven
federally-sourced community character signals.
Expert testimony establishing the quantitative basis for community-of-interest
preservation under the composite index is provided in [expert report reference].
The composite index contains no electoral, racial, or partisan data.
The plan satisfies the [state] redistricting criterion of [cite specific
criterion language] by preserving communities with composite scores above
[threshold] within single districts.''
\end{quote}
